\section{Related Work}
\label{sec:related_work}

\paragraph{Unsupervised graph domain adaptation.}
Early graph-transfer systems learn domain-invariant node representations with
adversarial alignment or shared encoders~\citep{zhang2019dane,wu2020unsupervised,dai2022graph}.
Later work makes the shift model more graph-specific: StruRW corrects
conditional structure, GRADE measures subtree discrepancy, A2GNN revisits
propagation, and PairAlign separates label and conditional-structure effects
~\citep{liu2023structural,wu2023non,liu2024rethinking,liu2024pairwisealignmentimprovesgraph}.
Recent methods additionally target smoothing, attribute-driven mismatch,
disentangled spectral shift, homophily, or composite distributions
~\citep{chen2025smoothness,fang2025attribute,yang2025disentangled,
pmlr-v267-fang25g,chen2026learningadaptivedistributionalignment}.
GDABench consolidates 16 algorithms and 74 transfers for systematic
comparison~\citep{liu2024revisitingbenchmarkingunderstandingunsupervised}.
These methods optimize the adapted predictor; \method instead holds all
candidate predictors fixed and studies how to choose among them without target
labels.

\paragraph{Label-free model selection under domain shift.}
Unsupervised domain adaptation has developed proxy criteria based on target
entropy, feature geometry, neighborhood density, transferability, and
cross-model agreement. DEV estimates target risk by embedding adapted features
into importance-weighted validation~\citep{you2019dev}; SND scores the density
of target features in a model-induced neighborhood graph
~\citep{saito2021tunerightwayunsupervised}; and Transfer Score combines
transferability and discriminability without target labels
~\citep{yang2023transferscore}. Agreement-based criteria exploit empirical
relations between classifier agreement and accuracy under shift
~\citep{jiang2021disagreement,baek2022agreement}, while recent work combines
multiple unsupervised signals for practical algorithm selection
~\citep{xiong2026practicalalgorithmselectionunsupervised}.
Such criteria are relevant baselines, but they treat the target examples as a
generic sample set and do not use the observed target adjacency to localize
agreement errors by graph frequency. \benchmark further evaluates one joint
candidate bank spanning algorithms, hyperparameters, seeds, and checkpoints
under an auditable sealed-label protocol.

\paragraph{Unsupervised classifier accuracy estimation.}
A complementary literature estimates classifier accuracy from unlabeled
prediction agreement. Spectral meta-learning ranks predictors when their
off-diagonal covariance has rank-one structure~\citep{parisi2014ranking}, and
moment-based estimators recover classifier accuracies under related low-rank or
conditional-independence assumptions~\citep{pmlr-v38-jaffe15}. UGDA trajectory
banks violate the spirit of these
assumptions because checkpoints from one run and methods sharing a GNN backbone
produce strongly correlated errors. Our disagreement identity makes this
dependence explicit through a band-wise error covariance term, and our recovery
bound quantifies how covariance misspecification propagates to risk error.

\paragraph{Graph spectral representations.}
Graph spectral methods decompose signals with functions of the graph
Laplacian~\citep{shuman2013emerging}. Spectral graph wavelets provide localized
multiscale atoms~\citep{hammond2011wavelets}; tight graph frames preserve signal
energy~\citep{shuman2013spectrumadaptedtightgraphwavelet}; and Chebyshev
polynomials enable localized filtering without eigendecomposition
~\citep{defferrard2016convolutional}.
Spectral ideas have also been used inside adaptation algorithms to align graph
bases, regularize transfer, or augment graph signals
~\citep{xiao2023spa,pang2023sa,yang2025disentangled}. In contrast, \method uses a tight
spectral frame after training: it decomposes candidate prediction errors,
recovers individual target risk from pairwise disagreement, and calibrates the
otherwise hidden covariance separately in each band. This distinction is
important because a tight frame alone cannot improve a linear selector; the
gain must arise from band-conditioned calibration, reliability, constraints,
or uncertainty.
